% Reviewed source SHA256 (UTF-8/LF): e6919a807c600c3774f722955e2aeeeb6cee649ad3abfb3428acc3003f9bce2e
\documentclass[11pt]{article}
\usepackage[a4paper,margin=27mm]{geometry}
\usepackage[T1]{fontenc}
\usepackage{lmodern,amsmath,amssymb,amsthm,mathtools,microtype}
\usepackage[hidelinks]{hyperref}
\newtheorem{theorem}{Theorem}\newtheorem{lemma}[theorem]{Lemma}
\newtheorem{corollary}[theorem]{Corollary}\newtheorem{proposition}[theorem]{Proposition}
\newcommand{\R}{\mathbb R}\newcommand{\sm}{\sigma_{\min}}
\title{A sharp weighted planar lemma and a subframe singular-value bound}
\author{Matthew J. Colbrook\\\small Department of Applied Mathematics and Theoretical Physics\\\small University of Cambridge, Cambridge, United Kingdom\\\small\href{mailto:m.colbrook@damtp.cam.ac.uk}{m.colbrook@damtp.cam.ac.uk}}\date{11 September 2026}
\usepackage{xurl}
\setlength{\emergencystretch}{3em}
\hypersetup{pdfauthor={Matthew J. Colbrook}}
\begin{document}\maketitle

\noindent\textbf{Independently reviewed partial result.} This is a polynomial auxiliary bound. The universal exponential conjecture remains open; novelty of the polynomial order is not claimed.
The dated \href{https://github.com/MColbrook/OpenProblemsInNLA/blob/codex/colbrook-frames-submissions/references/colbrook-frames-2026-09-11/verification/reviews/FR-04-review.md}{independent agent review} records the subsequent checks and supersedes original draft remarks about pending review. This is not external human peer review or formal certification.
\par\medskip
% BEGIN REVIEWED BODY

\begin{abstract}
For $M$ vectors in the real plane with prescribed total squared norm $E$, the smallest least singular value of a pair is at most $\sqrt{2E/M}\sin(\pi/(2M))$. Equality is attained by equal-length vectors on equally spaced projective lines. Combining this sharp weighted lemma with determinant-weighted projection gives an explicit bound for the smallest least singular value of a square row submatrix of any real rectangular matrix. At the minimally redundant phase-retrieval size $(2n-1)\times n$, it has order $n^{-3/2}$ times the largest row norm. The exponential bound requested in FR-04 is not obtained. The polynomial order was already present in the literature; this note supplies a self-contained weighted argument, not a resolution of FR-04.
\end{abstract}

\section{The weighted planar problem}
For row vectors $b_1,\ldots,b_M\in\R^2$, put
\[
 E=\sum_{i=1}^M\|b_i\|_2^2,\qquad
 t=\min_{i<j}\sm\begin{pmatrix}b_i\\b_j\end{pmatrix}.
\]

\begin{theorem}\label{thm:plane}
For every $M\ge2$,
\[
 t\le\sqrt{\frac{2E}{M}}\sin\frac{\pi}{2M}.
\]
The constant is sharp. If $E>0$, equality holds for vectors of norm $\sqrt{E/M}$ whose unoriented lines are equally spaced in the real projective line.
\end{theorem}
\begin{proof}
The case $t=0$ is immediate. Assume $t>0$ and write $q_i=\|b_i\|^2/t^2\ge1$. The vectors lie on distinct projective lines. Order the lines cyclically, with consecutive angle gaps $\theta_i\in(0,\pi)$ satisfying $\sum_i\theta_i=\pi$; indices below are cyclic.

Every two-row Gram matrix minus $t^2 I$ is positive semidefinite. Its determinant inequality implies
\[
 |\cos\theta_i|\le\sqrt{1-q_i^{-1}}\sqrt{1-q_{i+1}^{-1}}.
\]
Define $g(z)=\arccos(e^z)$ for $z<0$. Since
\[
 g''(z)=-\frac{e^z}{(1-e^{2z})^{3/2}}<0,
\]
concavity, with endpoint values interpreted by limits, gives
\begin{align*}
 \theta_i
 &\ge\arccos\left(\sqrt{1-q_i^{-1}}\sqrt{1-q_{i+1}^{-1}}\right)\\
 &\ge\frac12\arccos(1-q_i^{-1})+
       \frac12\arccos(1-q_{i+1}^{-1}).
\end{align*}
The function $f(q)=\arccos(1-q^{-1})$ on $q\ge1$ is convex, because for $q>1$,
\[
 f''(q)=\frac{3q-1}{q^2(2q-1)^{3/2}}>0.
\]
Consequently
\[
 \pi=\sum_i\theta_i\ge\sum_i f(q_i)
 \ge M f\left(\frac{E}{Mt^2}\right).
\]
As $f$ is decreasing, this is equivalent to
\[
 \frac{Mt^2}{E}\le1-\cos\frac\pi M
 =2\sin^2\frac\pi{2M},
\]
which proves the bound.

For equal-length vectors on equally spaced projective lines, the smallest angle between two lines is $\pi/M$. A pair with equal row norm $r$ and projective separation $\alpha$ has least singular value $r\sqrt{1-|\cos\alpha|}$. Thus its smallest value over all pairs is $\sqrt{2E/M}\sin(\pi/(2M))$, proving sharpness.
\end{proof}

The unit-maximum-row-norm planar extremum $\sqrt2 L\sin(\pi/(2M))$ is discussed by Liu~\cite{Liu}. Theorem~\ref{thm:plane} uses total squared norm instead and is no weaker, since $E/M\le L^2$. An exhaustive historical audit of this weighted formulation is not claimed.

\section{Determinant-weighted residual energy}
Let $A$ be a real $m\times n$ matrix, with rows $a_i$, $m\ge n\ge2$, and write $E=\|A\|_F^2$. For a row set $S$, let $P_S$ be the orthogonal projector onto the complement of the row span of $A_S$.

\begin{lemma}\label{lem:energy}
If $A$ has rank $n$, some linearly independent set $S$ of $n-2$ rows satisfies
\[
 \sum_{i=1}^m\|P_Sa_i\|_2^2\le\frac{2E}{n}.
\]
For $n=2$, use the empty set and $P_S=I$.
\end{lemma}
\begin{proof}
Let $k=n-2$ and give each $k$-element row set weight
$w_S=\det(A_S A_S^T)$. The weights are nonnegative and their sum is positive. If $w_S>0$, the Gram determinant formula gives
\[
 w_S\|P_Sa_i\|^2=\det(A_{S\cup\{i\}}A_{S\cup\{i\}}^T)
 \qquad(i\notin S).
\]
If $w_S=0$, both sides are zero: adding a single row to a dependent $k$-set cannot make a $(k+1)$-set independent. Summing and counting each $(k+1)$-set $k+1$ times yields
\[
 \frac{\sum_S w_S\sum_i\|P_Sa_i\|^2}{\sum_Sw_S}
 =(k+1)\frac{e_{k+1}(\lambda)}{e_k(\lambda)},
\]
where $\lambda_1,\ldots,\lambda_n>0$ are the eigenvalues of $A^T A$ and $e_j$ denotes the $j$th elementary symmetric polynomial. The identification of the determinant sums with $e_j(\lambda)$ follows from Cauchy--Binet.

Newton's inequalities imply that the ratios of consecutive normalized elementary symmetric polynomials decrease, so
\[
 \frac{e_{n-1}(\lambda)/n}{e_{n-2}(\lambda)/\binom n2}
 \le\frac{e_1(\lambda)}{n}=\frac En.
\]
Multiplication by $n-1$ gives a weighted mean at most $2E/n$, and at least one positive-weight set is no larger than that mean. The case $n=2$ is also the direct identity $E=2E/n$.
\end{proof}

\begin{theorem}\label{thm:general}
For every real $m\times n$ matrix, $m\ge n\ge2$, put
\[
 \omega_n(A)=\min_{|T|=n}\sm(A_T),\qquad M=m-n+2.
\]
Then
\[
 \boxed{\quad
 \omega_n(A)\le\sqrt{\frac{4\|A\|_F^2}{nM}}\,
                       \sin\frac\pi{2M}.
 \quad}
\]
\end{theorem}
\begin{proof}
If $A$ has rank less than $n$, the left side is zero. Otherwise choose $S$ from Lemma~\ref{lem:energy}. Its orthogonal complement is a plane. The $M$ rows not in $S$, projected to this plane, have total squared norm at most $2E/n$. Apply Theorem~\ref{thm:plane} to obtain a pair $i,j\notin S$ and a unit vector $x$ in that plane such that
\[
 |a_i x|^2+|a_j x|^2\le\frac{4E}{nM}\sin^2\frac\pi{2M}.
\]
All rows in $S$ annihilate $x$, so the same upper bound holds for $\|A_{S\cup\{i,j\}}x\|^2$. Taking a least singular value proves the assertion.
\end{proof}

\begin{corollary}\label{cor:critical}
For a real $(2n-1)\times n$ matrix with largest row norm $L$,
\[
 \omega_n(A)\le
 2L\sqrt{\frac{2n-1}{n(n+1)}}\,
 \sin\frac\pi{2(n+1)}
 \sim \pi\sqrt2\,L\,n^{-3/2}
\]
as an asymptotic expression for the upper bound.
\end{corollary}
\begin{proof}
Substitute $m=2n-1$ and $\|A\|_F^2\le(2n-1)L^2$ in Theorem~\ref{thm:general}.
\end{proof}

\section{The unresolved exponential step}
FR-04 asks whether universal constants $C>0$ and $0<\beta<1$ bound the same quantity by $CL\beta^n$ for every minimally redundant full-spark matrix~\cite{Repo}. Corollary~\ref{cor:critical} does not provide such constants. A polynomial upper bound cannot be substituted for an exponential one. The order $n^{-3/2}$ already appears in the discussion by Balan and Wang~\cite{BW}, so the order itself must not be presented as a new solution.

The argument isolates two sources of loss. Determinant-weighted projection controls a mean residual energy, not the distribution of the least singular values across exponentially many bases. The last two-dimensional step controls only the $M$ directions in a single projected plane. Neither step supplies exponential amplification. Small determinants by themselves do not repair the gap: the determinant is the product of all singular values, and without lower bounds for the other singular values it need not force an exponentially small least singular value.

The accompanying script checks the derivative identities, the determinant-weighted energy identity on exact rational examples, the sharp equally spaced planar examples, and numerical instances of the general bound. These checks are supplementary to the proof above.

\begin{thebibliography}{9}
\bibitem{Repo} OpenProblemsInNLA, \emph{FR-04: frames and matrix designs}. \url{https://github.com/ajt60gaibb/OpenProblemsInNLA/tree/main/frames-and-matrix-designs/FR-04}.
\bibitem{BW} R. Balan and Y. Wang, \emph{Invertibility and robustness of phaseless reconstruction}, Applied and Computational Harmonic Analysis 38 (2015), 469--488. \url{https://www.cscamm.umd.edu/publications/ACHA2015paper_CS-16-05.pdf}.
\bibitem{Liu} \emph{On the minimal smallest singular value of subframes for signal recovery}, Operators and Matrices 18 (2024), 273--293. \url{https://files.ele-math.com/articles/oam-18-17.pdf}.
\end{thebibliography}
\end{document}

% END REVIEWED BODY
